% ============================================
% IX. CONCLUSION AND FUTURE WORK
% ============================================
\section{Conclusion and Future Work}
\label{sec:conclusion}

In this paper, we presented a systematic, empirical analysis of architectural patterns for high-throughput blockchain-based trading systems, grounded in the Horizon Globex trading platform deployed by our industrial partner, Horizon Globex Ireland, for live regulated trading operations. Our proposed symbol-sharded, lock-free architecture achieves 31.19~tx/sec, representing a \textbf{2.8$\times$ improvement} over conventional asynchronous multi-threaded designs. The full ablation from a deliberately naive synchronous baseline demonstrates a cumulative 125-fold gain, though the more salient findings emerge from the incremental transitions between architectures. Our ablation study revealed a clear narrative of migrating bottlenecks:

First, we demonstrated that naive multi-threading offers limited benefit (a 36\% throughput gain from Architecture~2 to~3), as performance is ultimately serialized by the blockchain's single-wallet nonce requirement. Second, and more critically, we showed that resolving this on-chain bottleneck with a multi-wallet strategy (Architecture~4) created a server-side crisis, increasing average latency by over 9$\times$ due to severe lock contention.

Our primary contribution is the resolution to this server-side bottleneck. By implementing a symbol-sharded, lock-free model (Architecture~5), we not only increased throughput by a further 29\% but also reduced the catastrophic latency by 44\%. This final architecture successfully shifted the bottleneck away from the submission layer, achieving a throughput of 31.19~tx/sec. This work underscores that for high-performance blockchain systems, a holistic approach that co-designs on-chain strategies with off-chain data structures is essential to achieving optimal performance. Importantly, because our submission layer can sustain rates well in excess of the blockchain's gas-based theoretical capacity (~17.93~tx/sec), it scales linearly with increases in underlying blockchain capacity, whether through Layer~2 solutions, alternative consensus mechanisms, or network upgrades, making it suitable for deployment on higher-throughput networks without fundamental redesign.

\subsection{Future Work}

While our proposed architecture successfully addresses the challenge of raw throughput, it opens several promising avenues for future research that are critical for real-world deployment in enterprise financial systems.

\textbf{Sustained Load Testing and Saturation Analysis.} Our current experiments used a fixed workload of 1,000 transactions, which was sufficient to demonstrate the relative performance of different architectural patterns and to identify migrating bottlenecks. However, this workload did not sustain transaction submission long enough to reach steady-state saturation of the blockchain's gas-based theoretical capacity (~17.93~tx/sec based on our network's 10,485,760 gas block limit and 146,194 gas per transaction). Future work should conduct extended load tests that continuously submit transactions at rates exceeding this theoretical ceiling to empirically characterize: (a) the sustained throughput under saturation, (b) mempool behavior and backpressure dynamics, (c) latency distribution under sustained load, and (d) the system's behavior when the submission rate exceeds the blockchain's confirmation capacity for extended periods. Such testing would provide a more complete picture of the architecture's production readiness. Notably, newer deployments of private Ethereum networks increasingly support substantially higher gas block limits, which may effectively remove any meaningful theoretical throughput ceiling. Investigating how our architecture scales when the blockchain layer is no longer the binding constraint presents an interesting avenue for future research.

\textbf{MEV Mitigation and Fair Ordering.} Our current model optimizes for speed but does not explicitly guarantee fair ordering of transactions within a block. This leaves the system potentially vulnerable to Maximal Extractable Value (MEV) extraction~\cite{qin2022mev}, even in a permissioned setting where validators are known. Future work should explore the integration of fair ordering protocols, such as commit-reveal schemes or threshold encryption of the mempool, to ensure that validators cannot front-run or sandwich trades~\cite{grimmelmann2024mev}. Analyzing the trade-off between fairness and throughput would be a valuable contribution.

\textbf{Cross-Shard Atomic Trading.} The symbol-sharded model excels at processing trades for individual symbols in parallel. However, it does not natively support atomic swaps or complex trades that involve multiple symbols simultaneously. Executing such trades would require a cross-shard transaction protocol. We propose investigating the implementation of a two-phase commit (2PC) or optimistic cross-shard transaction mechanism to enable atomic settlement across different symbol shards, a significant challenge in distributed ledger systems~\cite{zhang2023frontrunning,harris2002crossshard}.

\textbf{Generalizability to Other DLT Platforms.} Our experiments were conducted on a private, IBFT~2.0-based Ethereum network. To validate the generalizability of our architectural principles, future work should include deploying and benchmarking the same architecture on other leading permissioned DLTs, most notably Hyperledger Fabric. Given Fabric's different architectural and consensus model, a comparative analysis would provide crucial insights into how platform choice interacts with application-level design for high-performance workloads~\cite{ucbas2023performance,baliga2018hyperledger}.

\textbf{Dynamic Resource and Workload Management.} The current system assumes a relatively static set of trading symbols and workloads. A more robust, production-ready system would need to dynamically adapt to changing market conditions, such as sudden spikes in volume for a particular symbol or the introduction of new trading pairs. Future research could focus on developing adaptive resource allocation algorithms that can dynamically scale worker threads or rebalance sharding strategies in real-time.

\textbf{Enhanced Privacy and Confidentiality.} While our system operates on a permissioned network, all settled trades are recorded transparently on the ledger. For many enterprise use cases, confidentiality of trade data is a strict requirement. Future iterations should explore the integration of privacy-enhancing technologies, such as zero-knowledge proofs (e.g., zk-SNARKs) or trusted execution environments (TEEs), to enable confidential transactions without sacrificing auditability~\cite{bensasson2014zerocash}.
